\fichatitulo{55 · METODOLOGIA}{Livro · 2002}{Projetos de pesquisa}
{\small GIL, Antônio Carlos. \textit{Como elaborar projetos de pesquisa}. 4. ed. São Paulo: Atlas, 2002. \href{https://files.cercomp.ufg.br/weby/up/150/o/Anexo_C1_como_elaborar_projeto_de_pesquisa_-_antonio_carlos_gil.pdf}{PDF consultado}.\par}

\campo{Problema e objetivo}
Como estruturar um projeto de pesquisa claro, viável e coerente com o problema, os objetivos, os métodos e as etapas de execução.

\campo{Principal contribuição · síntese interpretativa}
Apresenta orientações para formular problemas e hipóteses, classificar pesquisas, realizar pesquisa bibliográfica, elaborar fichamentos e redigir projetos. O texto liga a qualidade do projeto à delimitação do problema e à adequação dos procedimentos.

\campo{Método / natureza da evidência}
Livro metodológico de orientação e sistematização; não é estudo empírico nem valida experimentalmente um desenho específico.

\campo{Resultado ou argumento principal}
Defende que o projeto explicite problema, objetivos, justificativa, metodologia, recursos e cronograma. A pesquisa bibliográfica inclui identificação, localização, compilação, fichamento, análise, interpretação e redação.

\campo{Excertos-chave · seleção literal}
\begin{itemize}
\excerto{Problema}{O problema deve ser formulado como pergunta}{cap. 2, p. 27.}
\excerto{Projeto}{Quais os elementos de um projeto de pesquisa?}{cap. 1, p. 20.}
\end{itemize}

\campo{Limitações e alcance}
\textbf{Fonte:} a versão consultada é a 4ª edição de 2002, não a 7ª edição de 2022 registrada originalmente. \textbf{Análise da ficha:} serve como orientação metodológica geral; não deve ser usada para sustentar decisões específicas de avaliação de RAG.

\campo{Aproveitamento na dissertação · proposta}
\textbf{Destino:} metodologia e planejamento. \textbf{Uso:} sustentar a explicitação de natureza, objetivos, etapas e coerência entre problema e procedimentos.

\registro{Fonte consultada: PDF local, 176 páginas; capítulos sobre problema, hipóteses, pesquisa bibliográfica, fichamento e redação de projeto. Versão efetivamente identificada: 4. ed., 2002. Chave: \texttt{gil2022projetosPesquisa}.}
